\section{\texorpdfstring{第~\ref{chap:curvature}~章}{曲率}}\label{sec:solutions-curvature-chapter}
\subsection{\texorpdfstring{第~\ref{sec:curvature}~节}{曲率}}\label{subsec:solutions-curvature}

\begin{description}[font=\normalfont,leftmargin=2em,labelwidth=1.5em]
\RGsolution{ex:tensoriality-of-curvature}{tensoriality-of-curvature}
\begin{enumerate}[label=(\alph*),ref=\alph*,leftmargin=*]
\item
\begin{enumerate}[label=(\roman*),ref=\roman*,leftmargin=*]
\item\label{sol:curvature-linearity-first-slot} 只需证明 $R$ 对所有 $X,X_1,X_2,Y,Z\in\mathfrak X(M)$ 和 $f\in C^\infty(M)$ 满足
\[
  R(X_1+X_2,Y)Z=R(X_1,Y)Z+R(X_2,Y)Z
\]
和
\[
  R(fX,Y)Z=fR(X,Y)Z.
\]
第一个恒等式显然成立. 为证明第二个, 注意 $[fX,Y]=f[X,Y]-(Y\cdot f)X$, 因而
\begin{align*}
  R(fX,Y)Z&=\nabla_{fX}\nabla_YZ-\nabla_Y\nabla_{fX}Z-\nabla_{[fX,Y]}Z\\
  &=f\nabla_X\nabla_YZ-\nabla_Y(f\nabla_XZ)-\nabla_{f[X,Y]-(Y\cdot f)X}Z\\
  &=f\nabla_X\nabla_YZ-f\nabla_Y\nabla_XZ-(Y\cdot f)\nabla_XZ\\
  &\quad-f\nabla_{[X,Y]}Z+(Y\cdot f)\nabla_XZ\\
  &=f(\nabla_X\nabla_YZ-\nabla_Y\nabla_XZ-\nabla_{[X,Y]}Z)=fR(X,Y)Z.
\end{align*}
\item 与 (\ref{sol:curvature-linearity-first-slot}) 类似.
\item 同样, 只需证明 $R$ 对所有 $X,Y,Z_1,Z_2,Z\in\mathfrak X(M)$ 和 $f\in C^\infty(M)$ 满足
\[
  R(X,Y)(Z_1+Z_2)=R(X,Y)Z_1+R(X,Y)Z_2
\]
和
\[
  R(X,Y)(fZ)=fR(X,Y)Z.
\]
第一个恒等式显然成立. 第二个由下式得出:
\begin{align*}
  R(X,Y)(fZ)&=\nabla_X\nabla_Y(fZ)-\nabla_Y\nabla_X(fZ)-\nabla_{[X,Y]}(fZ)\\
  &=\nabla_X(f\nabla_YZ)+\nabla_X((Y\cdot f)Z)-\nabla_Y(f\nabla_XZ)\\
  &\quad-\nabla_Y((X\cdot f)Z)-f\nabla_{[X,Y]}Z-([X,Y]\cdot f)Z\\
  &=f\nabla_X\nabla_YZ+(X\cdot f)\nabla_YZ+(Y\cdot f)\nabla_XZ\\
  &\quad+(X\cdot(Y\cdot f))Z-f\nabla_Y\nabla_XZ-(Y\cdot f)\nabla_XZ\\
  &\quad-(X\cdot f)\nabla_YZ-(Y\cdot(X\cdot f))Z-f\nabla_{[X,Y]}Z\\
  &\quad-([X,Y]\cdot f)Z\\
  &=f(\nabla_X\nabla_YZ-\nabla_Y\nabla_XZ-\nabla_{[X,Y]}Z)\\
  &\quad+(X\cdot(Y\cdot f)-Y\cdot(X\cdot f))Z-([X,Y]\cdot f)Z\\
  &=fR(X,Y)Z.
\end{align*}
\end{enumerate}
\item 若 $x:V\to\mathbb{R}^n$ 为 $p\in M$ 附近的局部坐标, 则
\begin{align*}
  R(X,Y)Z
  &=R\left(\sum_{i=1}^nX^i\frac{\partial}{\partial x^i},
            \sum_{j=1}^nY^j\frac{\partial}{\partial x^j}\right)
     \left(\sum_{k=1}^nZ^k\frac{\partial}{\partial x^k}\right)\\
  &=\sum_{i,j,k=1}^nX^iY^jZ^kR\left(\frac{\partial}{\partial x^i},\frac{\partial}{\partial x^j}\right)
       \frac{\partial}{\partial x^k}\\
  &=\sum_{i,j,k,l=1}^n X^iY^jZ^kR_{ijk}{}^l\frac{\partial}{\partial x^l}.
\end{align*}
因此
\[
  (R(X,Y)Z)_p=\sum_{i,j,k,j=1}^nX^i(p)Y^j(p)Z^k(p)R_{ijk}{}^l(p)
   \left(\frac{\partial}{\partial x^l}\right)_p
\]
只依赖于 $X_p,Y_p,Z_p$. 此外, 这种依赖是线性的, 因而 $R$ 定义了一个 $(3,1)$-张量.
\end{enumerate}

\RGsolution{ex:sectional-curvature-independent-of-basis}{sectional-curvature-independent-of-basis}
设 $(X_p^1,X_p^2)$ 和 $(Y_p^1,Y_p^2)$ 是 $X_p$ 和 $Y_p$ 在由这两个向量生成的平面的一组标准正交基下的分量. 则 $X_p,Y_p$ 张成的平行四边形的面积平方为
\begin{align*}
 \begin{vmatrix}X_p^1&Y_p^1\\X_p^2&Y_p^2\end{vmatrix}^{\!2}
  &=(X_p^1Y_p^2-X_p^2Y_p^1)^2\\
  &=(X_p^1)^2(Y_p^2)^2+(X_p^2)^2(Y_p^1)^2-2X_p^1X_p^2Y_p^1Y_p^2.
\end{align*}
另一方面, 有
\begin{align*}
  &\lVert X_p\rVert^2\lVert Y_p\rVert^2-\langle X_p,Y_p\rangle^2\\
  &\quad=\bigl((X_p^1)^2+(X_p^2)^2\bigr)\bigl((Y_p^1)^2+(Y_p^2)^2\bigr)
      -\bigl(X_p^1Y_p^1+X_p^2Y_p^2\bigr)^2\\
  &\quad=(X_p^1)^2(Y_p^2)^2+(X_p^2)^2(Y_p^1)^2-2X_p^1X_p^2Y_p^1Y_p^2.
\end{align*}
若 $Z_p,W_p$ 是 $X_p,Y_p$ 生成的平面的另一组基, 则它们在该标准正交基下的分量满足
\[
  \begin{pmatrix}Z_p^1&W_p^1\\Z_p^2&W_p^2\end{pmatrix}
  =\begin{pmatrix}X_p^1&Y_p^1\\X_p^2&Y_p^2\end{pmatrix}S,
\]
其中 $S$ 是换基矩阵, 因而
\[
  \begin{vmatrix}Z_p^1&W_p^1\\Z_p^2&W_p^2\end{vmatrix}^{\!2}
  =(\det S)^2\begin{vmatrix}X_p^1&Y_p^1\\X_p^2&Y_p^2\end{vmatrix}^{\!2}.
\]
最后,
\begin{align*}
  R(Z_p,W_p,Z_p,W_p)
  &=R(S_{11}X_p+S_{21}Y_p,S_{12}X_p+S_{22}Y_p,\\
  &\qquad S_{11}X_p+S_{21}Y_p,S_{12}X_p+S_{22}Y_p)\\
  &=\bigl(S_{11}S_{22}S_{11}S_{22}-S_{11}S_{22}S_{21}S_{12}
       -S_{21}S_{12}S_{11}S_{22}\\
  &\qquad+S_{21}S_{12}S_{21}S_{12}\bigr)R(X_p,Y_p,X_p,Y_p)\\
  &=(S_{11}S_{22}-S_{12}S_{21})^2R(X_p,Y_p,X_p,Y_p)\\
  &=(\det S)^2R(X_p,Y_p,X_p,Y_p).
\end{align*}

\RGsolution{ex:rescaling-metric-and-curvature}{rescaling-metric-and-curvature}
\begin{enumerate}[label=(\alph*),ref=\alph*,leftmargin=*]
\item 首先注意, 相应的 Levi–Civita 联络 $\nabla^1$ 和 $\nabla^2$ 一致. 实际上, 由 Koszul 公式,
\begin{align*}
  2\langle\nabla_X^1Y,Z\rangle_1
  &=X\cdot\langle Y,Z\rangle_1+Y\cdot\langle X,Z\rangle_1-Z\cdot\langle X,Y\rangle_1\\
  &\quad-\langle[X,Z],Y\rangle_1-\langle[Y,Z],X\rangle_1+\langle[X,Y],Z\rangle_1\\
  &=2\rho\langle\nabla_X^2Y,Z\rangle_2,
\end{align*}
故对任意向量场 $X,Y,Z\in\mathfrak X(M)$, 有 $2\rho\langle\nabla_X^1Y,Z\rangle_2=2\rho\langle\nabla_X^2Y,Z\rangle_2$, 从而 $\nabla_X^1Y=\nabla_X^2Y$. 因此,
\begin{align*}
  R_1(X,Y,Z,W)
  &=\langle R_1(X,Y)Z,W\rangle_1=\langle R_2(X,Y)Z,W\rangle_1\\
  &=\rho\langle R_2(X,Y)Z,W\rangle_2=\rho R_2(X,Y,Z,W)
\end{align*}
所以对于 $M$ 的任一切空间的任意二维截面 $\Pi$, 有
\begin{align*}
  K_1(\Pi)&=-\frac{R_1(X,Y,X,Y)}{\lVert X\rVert_1^2\lVert Y\rVert_1^2-\langle X,Y\rangle_1^2}\\
  &=-\frac{\rho R_2(X,Y,X,Y)}
              {\rho^2(\lVert X\rVert_2^2\lVert Y\rVert_2^2-\langle X,Y\rangle_2^2)}
  =\rho^{-1}K_2(\Pi).
\end{align*}
\item 有
\[
  (R_1)_{ij}=\sum_{k=1}^n(R_1)_{kij}{}^k=\sum_{k=1}^n(R_2)_{kij}{}^k=(R_2)_{ij},
\]
因而 $\mathit{Ric}_1=\mathit{Ric}_2$.
\item 因为 $(g_1)_{ij}=\rho(g_2)_{ij}$ 蕴含 $(g_1)^{ij}=\rho^{-1}(g_2)^{ij}$, 所以
\[
  S_1(p)=\sum_{i,j=1}^n(R_1)_{ij}(g_1)^{ij}
  =\rho^{-1}\sum_{i,j=1}^n(R_2)_{ij}(g_2)^{ij}=\rho^{-1}S_2(p).
\]
\end{enumerate}
\end{description}
